%% SYNTHETIC TEST UNIT generated by writing/tools/make_outline_book.py from FINAL_OUTLINE_v2 -- not book text
\chapter{OPE：微扰证明、无质量 QCD、非微扰条件定理}\label{app:appM}
本单元为测试用的合成文本（由 FINAL\_OUTLINE\_v2 自动生成），陈述内容从略。

\section{第 1 节（测试）}\label{sec:appM-overview}

\begin{definition}[图、表观发散度、森林与 $R$ 运算]\label{def:appM-bphz}
（测试用合成陈述：内容从略。）
% uses: none
\end{definition}

\begin{lemma}[表观发散度的外线公式]\label{lem:appM-degree}
（测试用合成陈述：内容从略。）
\uses{def:appM-bphz}
\end{lemma}

\begin{theorem}[Zimmermann 收敛定理（顶点可加的超减除） \tagR]\label{thm:appM-zimmermann}
（$R_{\mathrm{PT}}$：在重整化微扰论中对红外正规化的部分子态严格。）
（测试用合成陈述：内容从略。）
（本命题只涉及部分子态。）
\uses{def:appM-bphz}
\end{theorem}

\begin{example}[非顶点可加的超减除会失败]\label{ex:appM-nonadditive}
（测试用合成陈述：内容从略。）
\uses{thm:appM-zimmermann, lem:appM-degree}
\end{example}

\begin{theorem}[Weinberg 定理（Hahn–Zimmermann 形式）]\label{thm:appM-weinberg}
（测试用合成陈述：内容从略。）
% uses: none
\end{theorem}

\begin{corollary}[渐近指数与子图发散度]\label{cor:appM-weinberg-degree}
（测试用合成陈述：内容从略。）
\uses{thm:appM-weinberg, def:appM-bphz}
\end{corollary}

\begin{remark}[森林公式不能逐项用 Weinberg 定理]\label{rem:appM-forest-asymptotics}
（测试用合成陈述：内容从略。）
% uses: none
\end{remark}

\section{第 2 节（测试）}\label{sec:appM-bphz}

\begin{definition}[正规乘积与冻结 $q$ 的双局域正规乘积]\label{def:appM-normal-product}
（测试用合成陈述：内容从略。）
\uses{def:appM-bphz}
\end{definition}

\begin{lemma}[收缩与桥]\label{lem:appM-contraction}
（测试用合成陈述：内容从略。）
\uses{def:appM-bphz}
\end{lemma}

\begin{lemma}[链展开恒等式]\label{lem:appM-chain-expansion}
（测试用合成陈述：内容从略。）
\uses{def:appM-bphz, lem:appM-contraction}
\end{lemma}

\begin{theorem}[Zimmermann 恒等式 \tagR]\label{thm:appM-zimmermann-identity}
（$R_{\mathrm{PT}}$：在重整化微扰论中对红外正规化的部分子态严格。）
（测试用合成陈述：内容从略。）
（本命题只涉及部分子态。）
\uses{def:appM-normal-product, lem:appM-chain-expansion, thm:appM-zimmermann}
\end{theorem}

\begin{corollary}[化到最小正规乘积 \tagR]\label{cor:appM-reduction}
（$R_{\mathrm{PT}}$：在重整化微扰论中对红外正规化的部分子态严格。）
（测试用合成陈述：内容从略。）
\uses{thm:appM-zimmermann-identity}
\end{corollary}

\begin{example}[$\varphi^4$ 中的 Zimmermann 恒等式]\label{ex:appM-callan-symanzik}
（测试用合成陈述：内容从略。）
\uses{thm:appM-zimmermann-identity}
\end{example}

\begin{lemma}[冻结 $q$ 时外动量分配的相容性]\label{lem:appM-routing}
（测试用合成陈述：内容从略。）
\uses{def:appM-normal-product, lem:appM-contraction}
\end{lemma}

\section{第 3 节（测试）}\label{sec:appM-chain}

\begin{lemma}[冻结 $q$ 正规乘积的良定义与链分解 \tagR]\label{lem:appM-integrability}
（$R_{\mathrm{PT}}$：在重整化微扰论中对红外正规化的部分子态严格。）
（测试用合成陈述：内容从略。）
\uses{def:appM-normal-product, lem:appM-chain-expansion, lem:appM-routing, thm:appM-zimmermann}
\end{lemma}

\begin{lemma}[接触项 \tagR]\label{lem:appM-contact}
（$R_{\mathrm{PT}}$：在重整化微扰论中对红外正规化的部分子态严格。）
（测试用合成陈述：内容从略。）
\uses{def:appM-normal-product, thm:appM-zimmermann}
\end{lemma}

\begin{theorem}[微扰论中的 OPE 恒等式（有质量模型） \tagR]\label{thm:appM-ope-identity}
（$R_{\mathrm{PT}}$：在重整化微扰论中对红外正规化的部分子态严格。）
（测试用合成陈述：内容从略。）
（本命题只涉及部分子态。）
\uses{lem:appM-integrability, lem:appM-contact, cor:appM-reduction, lem:appM-routing}
\end{theorem}

\begin{example}[树图 OPE 就是 Taylor 展开]\label{ex:appM-tree}
（测试用合成陈述：内容从略。）
\uses{thm:appM-ope-identity}
\end{example}

\begin{remark}[指数 $D-1-a$ 的启发式来源]\label{rem:appM-mechanism}
（测试用合成陈述：内容从略。）
% uses: none
\end{remark}

\begin{assumption}[冻结 $q$ 双局域正规乘积的动量空间余项估计 \tagA]\label{asm:appM-remainder}
（测试用合成陈述：内容从略。）
地位：测试状态句。
文献：测试。
被使用于：\cref{thm:appM-scaling}、\cref{cor:appM-coordinate}。
解除途径：测试解除途径。
\uses{def:appM-normal-product}
\end{assumption}
\begin{proofsketch}测试。\end{proofsketch}

\begin{lemma}[由齐次性、O(4) 与 Lagrange 恒等式得到的导数界]\label{lem:appM-derivative-bounds}
（测试用合成陈述：内容从略。）
\uses{def:appM-normal-product}
\end{lemma}

\section{第 4 节（测试）}\label{sec:appM-remainder}

\begin{definition}[精确投影、块与条件 (NE)]\label{def:appM-blocks}
（测试用合成陈述：内容从略。）
\uses{def:prelim-qcd, def:dis-twist}
\end{definition}

\begin{lemma}[条件 (NE) 的适用范围]\label{lem:appM-ne}
（测试用合成陈述：内容从略。）
\uses{def:appM-blocks}
\end{lemma}

\begin{lemma}[OPE 系数的唯一性 \tagR]\label{lem:appM-uniqueness}
（$R_{\mathrm{PT}}$：在重整化微扰论中对红外正规化的部分子态严格。）
（测试用合成陈述：内容从略。）
\uses{def:appM-normal-product, def:appM-blocks, lem:appM-ne}
\end{lemma}

\begin{theorem}[约化系数的标度与余项（有质量模型） \tagR]\label{thm:appM-scaling}
（$R_{\mathrm{cond}}$：在所列编号假设之下证明完整；结论的地位不高于其中最弱的假设。）
前提：\cref{asm:appM-remainder}。
（测试用合成陈述：内容从略。）
（本命题只涉及部分子态。）
% errtype: E-P
\uses{asm:appM-remainder, thm:appM-ope-identity, lem:appM-derivative-bounds, lem:appM-uniqueness}
\end{theorem}

\begin{example}[未约化系数违反标度（锐化值）]\label{ex:appM-unreduced}
（测试用合成陈述：内容从略。）
\uses{thm:appM-scaling, cor:appM-reduction}
\end{example}

\begin{lemma}[动量空间余项到坐标空间余项]\label{lem:appM-symbol}
（测试用合成陈述：内容从略。）
% uses: none
\end{lemma}

\begin{corollary}[坐标空间的 OPE 余项 \tagR]\label{cor:appM-coordinate}
（$R_{\mathrm{cond}}$：在所列编号假设之下证明完整；结论的地位不高于其中最弱的假设。）
前提：\cref{asm:appM-remainder}。
（测试用合成陈述：内容从略。）
（本命题只涉及部分子态。）
\uses{asm:appM-remainder, thm:appM-scaling, lem:appM-symbol, lem:appM-contact}
\end{corollary}

\section{第 5 节（测试）}\label{sec:appM-massless}

\begin{remark}[未识别的多项式项]\label{rem:appM-polynomial}
（测试用合成陈述：内容从略。）
% uses: none
\end{remark}

\begin{remark}[余项估计化归为零动量匹配估计；唯一性论证的循环性]\label{rem:appM-remainder-reduction}
（测试用合成陈述：内容从略。）
% uses: none
\end{remark}

\begin{remark}[例外动量与前向矩阵元]\label{rem:appM-exceptional}
（测试用合成陈述：内容从略。）
% uses: none
\end{remark}

\begin{proposition}[零动量扭度 2 算符插入在维数正规化中的存在性 \tagA]\label{prop:appM-zero-momentum-insertion}
（测试用合成陈述：内容从略。）
\uses{def:dis-twist}
\end{proposition}
\begin{proofsketch}测试。\end{proofsketch}
\paragraph{地位与缺口}测试：本陈述的地位由\cref{def:dis-twist}、\cref{rem:appM-exceptional}决定。

\begin{assumption}[前向在壳无质量部分子矩阵元的 OPE 余项 \tagA]\label{asm:appM-forward}
（测试用合成陈述：内容从略。）
地位：测试状态句。
文献：测试。
被使用于：\cref{prop:dis-ope-forward}。
解除途径：测试解除途径。
\uses{def:dis-compton}
\end{assumption}
\begin{proofsketch}测试。\end{proofsketch}

\begin{definition}[规范不变算符、BRST 恰当算符与运动方程算符]\label{def:appM-alien}
（测试用合成陈述：内容从略。）
\uses{def:prelim-brst}
\end{definition}

\begin{theorem}[定域 BRST 上同调]\label{thm:appM-local-cohomology}
（测试用合成陈述：内容从略。）
\uses{def:appM-alien, lem:prelim-brst-nilpotent}
\end{theorem}

\section{第 6 节（测试）}\label{sec:appM-brst}

\begin{lemma}[形式维数下逐圈的定域 BRST 上同调]\label{lem:appM-lbc-formal}
（测试用合成陈述：内容从略。）
\uses{def:appM-alien, thm:appM-local-cohomology}
\end{lemma}

\begin{remark}[Joglekar–Lee 混合定理的结构]\label{rem:appM-joglekar-lee}
（测试用合成陈述：内容从略。）
% uses: none
\end{remark}

\begin{lemma}[外来算符的物理矩阵元为零（非前向或红外正规化） \tagR]\label{lem:appM-alien-vanish}
（$R_{\mathrm{PT}}$：在重整化微扰论中对红外正规化的部分子态严格。）
（测试用合成陈述：内容从略。）
（本命题只涉及部分子态。）
\uses{def:appM-alien, thm:prelim-slavnov-taylor}
\end{lemma}

\begin{remark}[前向在壳胶子矩阵元的红外微妙之处]\label{rem:appM-forward-aliens}
（测试用合成陈述：内容从略。）
% uses: none
\end{remark}

\begin{assumption}[OPE 的 BRST 结构：带一个复合旁观者的欧氏展开 (Eε) \tagA]\label{asm:appM-brst-ope}
（测试用合成陈述：内容从略。）
地位：测试状态句。
文献：测试。
被使用于：\cref{prop:dis-ope-gi}、\cref{thm:appM-brst-reduction}。
解除途径：测试解除途径。
\uses{def:appM-alien}
\end{assumption}
\begin{proofsketch}测试。\end{proofsketch}

\begin{theorem}[OPE 系数的 BRST 约化 \tagR]\label{thm:appM-brst-reduction}
（$R_{\mathrm{cond}}$：在所列编号假设之下证明完整；结论的地位不高于其中最弱的假设。）
前提：\cref{asm:appM-brst-ope}。
（测试用合成陈述：内容从略。）
（本命题只涉及部分子态。）
\uses{asm:appM-brst-ope, def:appM-alien, lem:appM-lbc-formal, lem:appM-uniqueness, lem:appM-alien-vanish, prop:appJ-local-operators, lem:dis-twist-classification}
\end{theorem}

\begin{lemma}[前向 2-森林与 Symanzik 多项式 \tagR]\label{lem:appM-forward-forests}
（$R_{\mathrm{PT}}$：在重整化微扰论中对红外正规化的部分子态严格。）
（测试用合成陈述：内容从略。）
（本命题只涉及部分子态。）
\uses{def:dis-compton}
\end{lemma}

\section{第 7 节（测试）}\label{sec:appM-forward}

\begin{lemma}[Newton 多面体控制与 Taylor 级数]\label{lem:appM-newton}
（测试用合成陈述：内容从略。）
\uses{lem:appM-forward-forests}
\end{lemma}

\begin{lemma}[拟阵形式与维数正规化的可容许性]\label{lem:appM-matroid}
（测试用合成陈述：内容从略。）
\uses{lem:appM-forward-forests}
\end{lemma}

\begin{theorem}[前向在壳 DR 振幅的逐图展开定理 \tagR]\label{thm:appM-forward-expansion}
（$R_{\mathrm{PT}}$：在重整化微扰论中对红外正规化的部分子态严格。）
（测试用合成陈述：内容从略。）
（本命题只涉及部分子态。）
\uses{def:dis-compton, lem:appM-forward-forests, lem:appM-newton, lem:appM-matroid}
\end{theorem}

\begin{example}[前向三角图：精确的 Taylor 系数]\label{ex:appM-forward-triangle}
（测试用合成陈述：内容从略。）
\uses{thm:appM-forward-expansion}
\end{example}

\begin{example}[离壳极限与在壳 DR 值的非一致性]\label{ex:appM-forward-nonuniform}
（测试用合成陈述：内容从略。）
\uses{ex:appM-forward-triangle}
\end{example}

\begin{lemma}[D 维中 $T_2$ 的投影与前向振幅的全纯性 \tagR]\label{lem:appM-t2-projector}
（$R_{\mathrm{PT}}$：在重整化微扰论中对红外正规化的部分子态严格。）
（测试用合成陈述：内容从略。）
（本命题只涉及部分子态。）
\uses{thm:appM-forward-expansion, def:dis-compton}
\end{lemma}

\begin{corollary}[每阶的前向在壳部分子 OPE \tagA]\label{cor:appM-forward-ope}
（测试用合成陈述：内容从略。）
\uses{thm:appM-forward-expansion, lem:appM-matroid, lem:appM-t2-projector, lem:appM-ne, prop:dis-ope-euclid, thm:appM-brst-reduction, prop:appJ-local-operators, lem:appM-alien-vanish}
\end{corollary}
\begin{proofsketch}测试。\end{proofsketch}
\paragraph{地位与缺口}测试：本陈述的地位由\cref{asm:appM-forward}、\cref{thm:appM-forward-expansion}、\cref{lem:appM-matroid}、\cref{lem:appM-t2-projector}、\cref{lem:appM-ne}、\cref{prop:dis-ope-euclid}、\cref{thm:appM-brst-reduction}、\cref{prop:appJ-local-operators}、\cref{lem:appM-alien-vanish}、\cref{rem:appM-forward-aliens}、\cref{asm:appM-brst-ope}、\cref{prop:dis-ope-forward}、\cref{ex:appM-forward-oneloop}决定。

\section{第 8 节（测试）}\label{sec:appM-np}

\begin{assumption}[前向在壳胶子振幅的 Ward 恒等式 \tagA]\label{asm:appM-oswi}
（测试用合成陈述：内容从略。）
地位：测试状态句。
文献：测试。
被使用于：\cref{lem:appM-gluon-regularity}。
解除途径：测试解除途径。
\uses{def:dis-compton, def:prelim-brst}
\end{assumption}
\begin{proofsketch}测试。\end{proofsketch}

\begin{lemma}[由在壳 Ward 恒等式得到胶子靶在 $\omega=0$ 的正则性 \tagR]\label{lem:appM-gluon-regularity}
（$R_{\mathrm{cond}}$：在所列编号假设之下证明完整；结论的地位不高于其中最弱的假设。）
前提：\cref{asm:appM-oswi}。
（测试用合成陈述：内容从略。）
（本命题只涉及部分子态。）
\uses{asm:appM-oswi, thm:appM-forward-expansion, lem:appM-t2-projector}
\end{lemma}

\begin{corollary}[部分子前向振幅在割平面上的全纯性、不连续性与矩 \tagA]\label{cor:appM-cut-plane}
（测试用合成陈述：内容从略。）
\uses{thm:appM-forward-expansion, lem:appM-t2-projector, cor:appM-forward-ope}
\end{corollary}
\begin{proofsketch}测试。\end{proofsketch}
\paragraph{地位与缺口}测试：本陈述的地位由\cref{lem:appM-t2-projector}、\cref{cor:appM-forward-ope}、\cref{lem:appM-gluon-regularity}、\cref{thm:appM-forward-expansion}、\cref{prop:dis-compton-analytic}决定。

\begin{example}[单圈精确到 $\varepsilon$ 的检验]\label{ex:appM-forward-oneloop}
（测试用合成陈述：内容从略。）
\uses{prop:dis-ope-forward}
\end{example}

\begin{lemma}[前向管全纯性与与 $Q$ 无关的增长指数 \tagR]\label{lem:appM-forward-tube}
（$R_{\mathrm{cond}}$：在所列编号假设之下证明完整；结论的地位不高于其中最弱的假设。）
前提：\cref{asm:intro-wightman}。
（测试用合成陈述：内容从略。）
非微扰地位：\cref{lem:appM-forward-tube}、\cref{asm:intro-wightman}、\cref{def:dis-compton}。
\uses{asm:intro-wightman, def:dis-compton}
\end{lemma}

\begin{lemma}[$\hat P=0$ 点的不变量论证 \tagR]\label{lem:appM-hatP-zero}
（$R_{\mathrm{def}}$：由算符定义直接推出。）
（测试用合成陈述：内容从略。）
\uses{def:dis-compton}
\end{lemma}

\begin{lemma}[沿弯曲路径逼近的边界值]\label{lem:appM-curved-approach}
（测试用合成陈述：内容从略。）
% uses: none
\end{lemma}

\section{第 9 节（测试）}\label{sec:appM-status}

\begin{lemma}[不连续性的正性（$Q^2$ 抹平） \tagR]\label{lem:appM-positivity}
（$R_{\mathrm{cond}}$：在所列编号假设之下证明完整；结论的地位不高于其中最弱的假设。）
前提：\cref{asm:intro-wightman}。
（测试用合成陈述：内容从略。）
非微扰地位：\cref{lem:appM-positivity}、\cref{asm:intro-wightman}、\cref{lem:appM-curved-approach}、\cref{lem:appM-forward-tube}、\cref{prop:dis-W-structure}。
\uses{asm:intro-wightman, lem:appM-forward-tube, prop:dis-W-structure, lem:appM-curved-approach}
\end{lemma}

\begin{lemma}[几乎解析延拓与 Green 公式：不连续性测度的增长]\label{lem:appM-green}
（测试用合成陈述：内容从略。）
% uses: none
\end{lemma}

\begin{lemma}[RG 演化矩阵的范数界（质量混合的重标度） \tagR]\label{lem:appM-norm-bound}
（$R_{\mathrm{cond}}$：在所列编号假设之下证明完整；结论的地位不高于其中最弱的假设。）
前提：\cref{asm:intro-composite}、\cref{asm:intro-af-coeff}。
（测试用合成陈述：内容从略。）
非微扰地位：\cref{lem:appM-norm-bound}、\cref{asm:intro-composite}、\cref{asm:intro-af-coeff}。
\uses{asm:intro-composite, asm:intro-af-coeff}
\end{lemma}

\begin{remark}[实用 OPE、基的重排与重整子吸收]\label{rem:appM-svz}
（测试用合成陈述：内容从略。）
% uses: none
\end{remark}

\begin{remark}[支持性的严格结果]\label{rem:appM-supporting}
（测试用合成陈述：内容从略。）
% uses: none
\end{remark}

\begin{remark}[瞬子与渐近性]\label{rem:appM-instantons}
（测试用合成陈述：内容从略。）
% uses: none
\end{remark}

\begin{remark}[真空 OPE]\label{rem:appM-vacuum}
（测试用合成陈述：内容从略。）
% uses: none
\end{remark}
